% !TEX root = ../main.tex
\section{Graph-Based Search Algorithms}

We implemented Uniform-Cost Search (uninformed) and \astar{} (informed) (\dec{01}). Both are one
best-first \textbf{graph search} in \code{search(grid, start, goal, algorithm)}; they differ only in
the priority:
\begin{itemize}
  \item \textbf{UCS} expands the frontier node with the lowest path cost $g(n)$. With non-negative
    step costs, it is complete and optimal: when a node is removed from the frontier, its cheapest
    path is known.
  \item \textbf{\astar} expands the node with the lowest $f(n) = g(n) + h(n)$, where $h$ estimates
    the remaining cost. With a consistent $h$ it is complete and optimal, and it expands no node
    with $f > C^*$ (the optimal cost). This is why it usually expands fewer nodes than UCS.
\end{itemize}

Shared implementation details: a binary heap is the frontier. Ties are broken by insertion order,
so results are deterministic. A node is goal-tested when it is \emph{popped}, not when it is
generated, because testing on generation can return a non-optimal path. A \code{closed} set ensures
each cell is expanded at most once. A node counts as expanded when its successors are generated
(\dec{14}). \code{search} returns \code{(path, cost, nodes\_expanded)}, or
\code{(None, None, nodes\_expanded)} when no path exists; for start $=$ goal it returns
\code{([start], 0, 0)}. An unknown algorithm name raises \code{ValueError}.

\textbf{Heuristic choice.} For \astar{} we use Manhattan distance,
$h(n) = |row_n - row_{goal}| + |col_n - col_{goal}|$ (spec §5.5; \dec{02}). It is admissible:
each move is UP, DOWN, LEFT or RIGHT at cost 1, so one move reduces the Manhattan distance by at
most 1, and obstacles can only lengthen the real path. Hence $h(n)$ never exceeds the true remaining
cost. It is also consistent: for any move from $n$ to a neighbour $n'$, $h(n) \le 1 + h(n')$, and
$h(goal) = 0$. Because our \astar{} is a graph search with an explored set, consistency guarantees
that each cell is expanded first along a cheapest path. So \astar{} returns an optimal path with the
same cost as UCS, without reopening cells. Among admissible options it is the most informed:
Euclidean distance and $h = 0$ are never larger, so they would make \astar{} expand more nodes, and
$h = 0$ reduces \astar{} to UCS.

\begin{algorithm}[H]
\caption{Best-first graph search shared by UCS and \astar}
\begin{algorithmic}[1]
\Function{Search}{$grid, start, goal, algorithm$}
  \State $priority(n, g) \gets g + h(n, goal)$ \textbf{if} $algorithm = $ \astar{} \textbf{else} $g$ \Comment{UCS: $h \equiv 0$}
  \State $frontier \gets$ heap $\{(priority(start,0), start)\}$; \ $best\_g[start] \gets 0$
  \State $parent[start] \gets \textit{none}$; \ $closed \gets \emptyset$; \ $expanded \gets 0$
  \While{$frontier \neq \emptyset$}
    \State $n \gets$ pop lowest priority (FIFO among ties)
    \If{$n \in closed$} \textbf{continue} \Comment{stale duplicate entry}
    \EndIf
    \If{\Call{Goal-Test}{$n, goal$}}
      \State \Return $(\Call{Path}{parent, n},\ best\_g[n],\ expanded)$
    \EndIf
    \State $closed \gets closed \cup \{n\}$; \ $expanded \gets expanded + 1$
    \For{$(m, action, step) \in \Call{Get-Successors}{n, grid}$} \Comment{UP, DOWN, LEFT, RIGHT}
      \State $g \gets best\_g[n] + step$
      \If{$m \notin closed$ \textbf{and} $g < best\_g.\mathrm{get}(m, \infty)$}
        \State $best\_g[m] \gets g$; \ $parent[m] \gets n$
        \State push $(priority(m, g), m)$ onto $frontier$
      \EndIf
    \EndFor
  \EndWhile
  \State \Return $(\textit{None}, \textit{None}, expanded)$ \Comment{goal unreachable}
\EndFunction
\end{algorithmic}
\end{algorithm}

UCS is this procedure with $h \equiv 0$. \astar{} uses \code{heuristic(state, goal\_state)}.
